\section{Turing Machine --- Code Format}
%%% generated by the blueprint pipeline from TCSlib.Complexity.TuringMachine.CodeFormat
\subsection*{Turing Machine --- Code Format: Introduction}

Declarations of the Lean module \texttt{TCSlib.Complexity.TuringMachine.CodeFormat}.

\subsection*{Turing Machine --- Code Format: Results}

\begin{definition}[Bit encoding of a machine action]
\lean{Turing.actionBitsK}
\label{Turing.actionBitsK}
\leanok
\uses{Turing.Action, Turing.optBoolBits, Turing.optOptBoolBits, Turing.optStateBits, Turing.signBits}
For an action of a machine with $k$ work tapes, Boolean symbols, and states in a set of
size $n+1$, form a bit string by concatenating the encoding of the input-head movement,
then for each work tape the encodings of its optional write instruction and head
movement, followed by the output-symbol encoding and successor-state encoding.
\end{definition}

\begin{definition}[Reading a work-tape action code]
\lean{Turing.codeReadWork}
\label{Turing.codeReadWork}
\leanok
\uses{Turing.codeReadSign, Turing.codeReadWrite}
Given a bit string, attempt to decode first an optional write instruction, consisting of
an optional optional Boolean value, and then a head movement. If both decodings succeed,
return the decoded pair together with the remaining suffix of the bit string; otherwise,
return no result.
\end{definition}

\begin{definition}[Action Decoder]
\lean{Turing.codeReadActionK}
\label{Turing.codeReadActionK}
\leanok
\uses{Turing.Action, Turing.codeReadOutput, Turing.codeReadSign, Turing.codeReadState, Turing.codeReadVec, Turing.codeReadWork}
For natural numbers $k,n$ and a Boolean list $x$, the decoder attempts to read, in
order, an input-head movement, the actions on $k$ work tapes, an optional output symbol,
and a state among $n+1$ possible states. On success it returns the resulting action,
whose state set is $	ext{Fin}(n+1)$, together with the unconsumed Boolean list;
otherwise it returns no result.
\end{definition}

\begin{definition}[Work Action Skipper]
\lean{Turing.codeSkipWork}
\label{Turing.codeSkipWork}
\leanok
\uses{Turing.codeSkipPair}
For a Boolean list $x$, first apply the pair-skipping operation with predicate always
true, then apply it to the resulting list with predicate $a\lor\neg b$ on each pair
$(a,b)$. The result is the remaining list if both operations succeed, and no result
otherwise.
\end{definition}

\begin{definition}[Sequential code-skipping action]
\lean{Turing.codeSkipActionK}
\label{Turing.codeSkipActionK}
\leanok
\uses{Turing.codeSkipPair, Turing.codeSkipRepeat, Turing.codeSkipState, Turing.codeSkipWork}
For natural numbers $k$ and $n$ and a Boolean list $xs$, apply the code-skipping
operations in this order: skip a pair using the Boolean condition $a \lor \neg b$,
repeat the work-skipping operation $k$ times, skip another pair using the same
condition, and then skip a state encoding for $n+1$. The result is the remaining Boolean
list if every operation succeeds, and no result if any operation fails.
\end{definition}

\begin{lemma}[Decoding a work action after encoding]
\lean{Turing.codeReadWork_append}
\label{Turing.codeReadWork_append}
\leanok
\uses{Turing.codeReadSign_append, Turing.codeReadWork, Turing.codeReadWrite_append, Turing.optOptBoolBits, Turing.signBits}
\difficulty{3}
For every optional work-tape write $w$, head movement $m\in\{-1,0,1\}$, and bit string
$x$, the work-action reader applied to the encoding of $w$ followed by the encoding of
$m$, and then by $x$, succeeds and returns the pair $(w,m)$ and suffix $x$.
\end{lemma}

\begin{lemma}[Decoding an encoded transition action]
\lean{Turing.codeReadActionK_append}
\label{Turing.codeReadActionK_append}
\leanok
\uses{Turing.Action, Turing.actionBitsK, Turing.codeReadActionK, Turing.codeReadOutput_append, Turing.codeReadSign_append, Turing.codeReadState_append, Turing.codeReadVec, Turing.codeReadVec_append, Turing.codeReadWork, Turing.codeReadWork_append, Turing.optOptBoolBits, Turing.signBits}
\difficulty{4}
For every natural number $k$, every natural number $n$, every action of a machine with
$k$ work tapes and state set $\{0,\ldots,n\}$, and every Boolean string $x$, reading the
serialization of that action followed by $x$ returns the action and the suffix $x$.
\end{lemma}

\begin{lemma}[Soundness of work-action decoding]
\lean{Turing.codeReadWork_sound}
\label{Turing.codeReadWork_sound}
\leanok
\uses{Turing.codeReadSign_sound, Turing.codeReadWork, Turing.codeReadWrite_sound, Turing.optOptBoolBits, Turing.signBits}
\difficulty{4}
If the work-action reader succeeds on a bit string $x$, returning an optional work-tape
write $w$, a head movement $m\in\{-1,0,1\}$, and suffix $r$, then $x$ is the encoding of
$w$ followed by the encoding of $m$, followed by $r$.
\end{lemma}

\begin{lemma}[Action reader soundness]
\lean{Turing.codeReadActionK_sound}
\label{Turing.codeReadActionK_sound}
\leanok
\uses{Turing.Action, Turing.actionBitsK, Turing.codeReadActionK, Turing.codeReadOutput_sound, Turing.codeReadSign_sound, Turing.codeReadState_sound, Turing.codeReadVec_sound, Turing.codeReadWork, Turing.codeReadWork_sound, Turing.optOptBoolBits, Turing.signBits}
\difficulty{4}
For every pair of natural numbers $k,n$, bit string $x$, action $a$ of a $k$-work-tape
binary machine with states indexed by $\{0,\ldots,n\}$, and bit string $r$, if reading
$x$ as an action with $k$ work tapes and parameter $n$ succeeds with result $(a,r)$,
then $x$ is the serialization of $a$ followed by $r$.
\end{lemma}

\begin{lemma}[Erasure law for work-tape actions]
\lean{Turing.codeEraseWork}
\label{Turing.codeEraseWork}
\leanok
\uses{Turing.codeEraseSign, Turing.codeEraseWrite, Turing.codeErase_bind, Turing.codeReadWork, Turing.codeSkipWork}
\difficulty{4}
For every bit string $x$, discarding the decoded work-tape write and head movement from
the work-action reader gives exactly the result of a skipper that first consumes any
two-bit pair, then consumes a second pair unless it is $01$.
\end{lemma}

\begin{lemma}[Erasure law for action records]
\lean{Turing.codeEraseActionK}
\label{Turing.codeEraseActionK}
\leanok
\uses{Turing.codeEraseOutput, Turing.codeEraseSign, Turing.codeEraseState, Turing.codeEraseVec, Turing.codeEraseWork, Turing.codeErase_bind, Turing.codeReadActionK, Turing.codeReadVec, Turing.codeReadWork, Turing.codeSkipActionK, Turing.codeSkipRepeat, Turing.codeSkipWork}
\difficulty{5}
For all natural numbers $k,n$ and every bit string $x$, discarding the parsed action and
retaining only the remaining suffix from the reader for a $k$-work-tape action over
$n+1$ states gives exactly the suffix returned by the corresponding action skipper.
\end{lemma}

\begin{lemma}[Minimum length of a multi-tape action code]
\lean{Turing.codeActionK_length}
\label{Turing.codeActionK_length}
\leanok
\uses{Turing.Action, Turing.actionBitsK, Turing.codeFlatMap_length, Turing.optBoolBits, Turing.optOptBoolBits, Turing.optStateBits, Turing.signBits, Turing.unaryFin}
\difficulty{4}
For every number $k$ of work tapes, every $n\in\mathbb N$, and every action of a
$k$-tape machine with Boolean symbols and state set of size $n+1$, the action’s bit
encoding has length at least $4k+5$.
\end{lemma}

\begin{definition}[Two Work Tape Entries]
\lean{Turing.workPair}
\label{Turing.workPair}
\leanok
Given two optional Boolean symbols $w_0,w_1$, define a function on the two indices $0,1$
that returns $w_0$ at index $0$ and $w_1$ at index $1$.
\end{definition}

\begin{definition}[Two-Tape Action Encoding]
\lean{Turing.actionBits₂}
\leanok
\uses{Turing.Action, Turing.optBoolBits, Turing.optOptBoolBits, Turing.optStateBits, Turing.signBits}
For a natural number $n$ and an action of a two-work-tape machine with Boolean alphabet
and state set of size $n+1$, encode its data as a Boolean list: the input-head movement,
then the write option and head movement for each work tape in order, then the optional
output symbol, and finally the optional successor state.
\end{definition}

\begin{lemma}[One-tape action serialization]
\lean{Turing.actionBits_eq_actionBitsK}
\label{Turing.actionBits_eq_actionBitsK}
\leanok
\uses{Turing.Action, Turing.actionBits, Turing.actionBitsK}
\difficulty{1}
For every $n\in\mathbb{N}$ and every action of a one-work-tape binary machine with state
set of size $n+1$, the serialization formed by listing its input-head move, its
work-tape write and head move, its output symbol, and its successor state agrees with
the serialization formed by listing the same data through the general $k$-tape procedure
with $k=1$.
\end{lemma}

\begin{lemma}[Two-tape action serialization]
\lean{Turing.actionBits₂_eq_actionBitsK}
\leanok
\uses{Turing.Action, Turing.actionBitsK, Turing.actionBits₂}
\difficulty{1}
For every $n\in\mathbb{N}$ and every action of a two-work-tape binary machine with state
set of size $n+1$, the serialization formed by listing its input-head move, both
work-tape writes and head moves, its output symbol, and its successor state agrees with
the serialization formed by listing the same data through the general $k$-tape procedure
with $k=2$.
\end{lemma}

\begin{lemma}[Equivalence of two transition readers]
\lean{Turing.codeSkipAction_eq_codeSkipActionK}
\label{Turing.codeSkipAction_eq_codeSkipActionK}
\leanok
\uses{Turing.codeSkipAction, Turing.codeSkipActionK, Turing.codeSkipRepeat, Turing.codeSkipWork}
\difficulty{4}
For every natural number $n$ and Boolean string $xs$, the transition-record reader that
checks four consecutive two-bit fields and then a successor-state field returns the same
result as the reader that checks the first field, skips one work-field pair, checks the
next field, and then checks the successor-state field.
\end{lemma}

\begin{definition}[Machine Code Format]
\lean{Turing.CodeFormat}
\label{Turing.CodeFormat}
\leanok
\uses{Turing.pairEncode, Turing.unaryFin}
A code format specifies a type of machine codes, a type of transition tables for each
natural number $n$, a state-count function, an assembly operation taking $n$, an initial
state in $\text{Fin}(n+1)$, and a table, and fixed serializations for codes and tables.
It also provides table-reading and erased table-skipping operations, a natural-number
guard, and a fallback code. These data satisfy the following conditions: assembled
machines have state count $n$; their serialization is the pair encoding of the binary
representation of $n$ followed by the unary encoding of the initial state and the table
bits; every code is assembled from some such data; reading a table's bits followed by
any suffix recovers that table and suffix; every successful table read splits its input
into those table bits followed by the returned remainder; erasing the result of table
reading agrees with skipping; and every table's bit length is at least the guard times
$n+1$.
\end{definition}

\begin{definition}[Full Code Parser]
\lean{Turing.CodeFormat.parseFull}
\label{Turing.CodeFormat.parseFull}
\leanok
\uses{Turing.codeBitsNat, Turing.codeReadFin, Turing.pairDecode}
For a code format and a Boolean list $x$, first decode a pair consisting of a bit list
and a remainder, and interpret the bit list as a natural number $n$. Continue only if
the bit list is exactly the binary representation of $n$ and the remainder has length at
least the format's guard times $n+1$; then read an initial state among $n+1$ states and
a transition table for $n$. Return the assembled code and the final remainder if every
bit of that remainder is true; otherwise return no result.
\end{definition}

\begin{definition}[Code Decoder]
\lean{Turing.CodeFormat.decode}
\label{Turing.CodeFormat.decode}
\leanok
\uses{Turing.CodeFormat.parseFull}
For a code format and a Boolean list $x$, return the code produced by full parsing when
parsing succeeds; if parsing fails, return the format's fallback code.
\end{definition}

\begin{definition}[Code format scanner]
\lean{Turing.CodeFormat.scan}
\label{Turing.CodeFormat.scan}
\leanok
\uses{Turing.codeBitsNat, Turing.codeSkipFin, Turing.pairDecode}
For a code format $F$, the scanner takes a bit string, decodes it as a pair of a length
code and a remainder, and lets $n$ be the natural number represented by the length code.
It returns no result unless that code is exactly the bit representation of $n$ and the
remainder has at least $g(n+1)$ bits, where $g$ is the format’s guard constant.
Otherwise it removes the unary encoding of a state index in the range of $n+1$ states,
then removes the encoded transition table; it succeeds with the remaining suffix
precisely when every bit in that suffix is $1$.
\end{definition}

\begin{definition}[Canonical bit-string representative]
\lean{Turing.CodeFormat.canonical}
\label{Turing.CodeFormat.canonical}
\leanok
\uses{Turing.CodeFormat.scan}
For a code format $F$ and a bit string $x$, define its canonical representative as
follows: if scanning $x$ succeeds and returns a tail $t$, take the prefix of $x$
obtained by removing the final $|t|$ bits. If scanning fails, use the serialization of
the format’s fallback code.
\end{definition}

\begin{theorem}[Parsing a padded serialization]
\lean{Turing.CodeFormat.parseFull_serialize_pad}
\label{Turing.CodeFormat.parseFull_serialize_pad}
\leanok
\uses{Turing.CodeFormat.parseFull, Turing.codeBitsNat_bits, Turing.codeReadFin_append, Turing.pairDecode_pairEncode, Turing.pairEncode}
\difficulty{5}
For every code $M$ and every natural number $m$, parsing the serialization of $M$
followed by $m$ copies of the bit $1$ succeeds and returns $M$ together with those $m$
remaining bits.
\end{theorem}

\begin{theorem}[Decoding a padded serialization]
\lean{Turing.CodeFormat.decode_serialize_pad}
\label{Turing.CodeFormat.decode_serialize_pad}
\leanok
\uses{Turing.CodeFormat.decode, Turing.CodeFormat.parseFull_serialize_pad}
\difficulty{1}
For every code $M$ and every natural number $m$, decoding the serialization of $M$
followed by $m$ copies of the bit $1$ returns $M$.
\end{theorem}

\begin{theorem}[Soundness of full parsing]
\lean{Turing.CodeFormat.parseFull_sound}
\label{Turing.CodeFormat.parseFull_sound}
\leanok
\uses{Turing.CodeFormat.parseFull, Turing.codeBitsNat, Turing.codeReadFin_sound, Turing.eq_pairEncode_of_pairDecode, Turing.pairEncode}
\difficulty{5}
For every bit string $x$, code $M$, and bit string $t$, if parsing $x$ succeeds with
output $M$ and remainder $t$, then $x$ is the serialization of $M$ followed by $t$.
\end{theorem}

\begin{theorem}[Scan returns the parsed remainder]
\lean{Turing.CodeFormat.scan_full}
\label{Turing.CodeFormat.scan_full}
\leanok
\uses{Turing.CodeFormat.parseFull, Turing.CodeFormat.scan, Turing.codeBitsNat, Turing.codeEraseFin, Turing.codeErase_bind}
\difficulty{5}
For every bit string $x$, scanning $x$ returns exactly the remainder obtained by parsing
$x$, when parsing succeeds; if parsing fails, scanning also fails.
\end{theorem}

\begin{theorem}[Canonical code is serialized decoding]
\lean{Turing.CodeFormat.canonical_eq}
\label{Turing.CodeFormat.canonical_eq}
\leanok
\uses{Turing.CodeFormat.canonical, Turing.CodeFormat.decode, Turing.CodeFormat.parseFull, Turing.CodeFormat.parseFull_sound, Turing.CodeFormat.scan_full}
\difficulty{4}
For every code format $F$ and every Boolean word $xs$, the canonical form of $xs$ equals
the serialization of the code obtained by decoding $xs$.
\end{theorem}

\begin{theorem}[Serialization length after decoding]
\lean{Turing.CodeFormat.serialize_decode_length_le}
\label{Turing.CodeFormat.serialize_decode_length_le}
\leanok
\uses{Turing.CodeFormat.canonical, Turing.CodeFormat.canonical_eq, Turing.CodeFormat.decode, Turing.CodeFormat.scan}
\difficulty{4}
For any code format $F$ and bit string $\alpha$, the length of the serialization of the
value obtained by decoding $\alpha$ is at most the larger of $|\alpha|$ and the
serialization length of the fallback value of $F$.
\end{theorem}

\begin{theorem}[Decoded state-count bound]
\lean{Turing.CodeFormat.numStates_decode_le}
\label{Turing.CodeFormat.numStates_decode_le}
\leanok
\uses{Turing.CodeFormat.decode, Turing.CodeFormat.serialize_decode_length_le, Turing.length_pairEncode}
\difficulty{4}
For every bit string $\alpha$, the guard constant of $F$ times one more than the number
of states of the code decoded from $\alpha$ is at most the maximum of the length of
$\alpha$ and the length of the serialization of $F$'s fallback code.
\end{theorem}

\begin{theorem}[Failed full parse yields fallback]
\lean{Turing.CodeFormat.decode_of_parseFull_none}
\label{Turing.CodeFormat.decode_of_parseFull_none}
\leanok
\uses{Turing.CodeFormat.decode, Turing.CodeFormat.parseFull}
\difficulty{1}
For every code format $F$ and Boolean word $\alpha$, if the full parser for $F$ fails on
$\alpha$, then decoding $\alpha$ by $F$ returns its fallback code.
\end{theorem}

\begin{theorem}[Insufficient suffix length forces fallback]
\lean{Turing.CodeFormat.decode_of_guard_lt}
\label{Turing.CodeFormat.decode_of_guard_lt}
\leanok
\uses{Turing.CodeFormat.decode, Turing.CodeFormat.parseFull, Turing.codeBitsNat, Turing.pairDecode}
\difficulty{1}
Let $F$ be a code format with guard parameter $g$, and let $\alpha$, $b$, and $r$ be
finite Boolean words. Suppose paired decoding of $\alpha$ succeeds with components $b$
and $r$, and let $m$ be the natural number represented by $b$ in
least-significant-bit-first binary. If $|r| < g(m+1)$, then decoding $\alpha$ by $F$
returns the fallback code of $F$.
\end{theorem}

\begin{definition}[Two-tape code table type]
\lean{Turing.CodeTable2}
\label{Turing.CodeTable2}
\leanok
\uses{Turing.Action}
For each natural number $n$, a two-tape code table with $n+1$ states assigns to every
state and every triple consisting of the possibly blank input symbol and the possibly
blank symbols under the two work-tape heads an action of a two-work-tape machine over
the Boolean alphabet, with state set of size $n+1$.
\end{definition}

\begin{definition}[Bit encoding of a two-tape table]
\lean{Turing.codeTable2Bits}
\label{Turing.codeTable2Bits}
\leanok
\uses{Turing.CodeTable2, Turing.actionBitsK}
Given a two-tape code table with $n+1$ states, encode its entries as a bit string by
listing the states in order and, for each state, listing the actions in lexicographic
product order over the three possible input symbols (blank, false, true) and the three
possible scanned symbols on each work tape (blank, false, true), encoding each action by
its action bit string.
\end{definition}

\begin{definition}[Two-tape table reader]
\lean{Turing.codeReadTable2}
\label{Turing.codeReadTable2}
\leanok
\uses{Turing.CodeTable2, Turing.codeReadActionK, Turing.codeReadSymbols, Turing.codeReadVec}
For a natural number $n$, the table reader takes a bit string and attempts to parse the
encoded actions for all $n+1$ states and all combinations of the possibly blank input
symbol and the two possibly blank work-tape symbols. On success it returns the resulting
two-tape code table together with the unconsumed suffix; on failure it returns no
result.
\end{definition}

\begin{definition}[Two-tape table skipper]
\lean{Turing.codeSkipTable2}
\label{Turing.codeSkipTable2}
\leanok
\uses{Turing.codeSkipActionK, Turing.codeSkipRepeat}
For a natural number $n$, the table skipper attempts to consume from a bit string the
encoded actions for a two-tape table with $n+1$ states, covering all three possibilities
for each of the input symbol and the two work-tape symbols at every state. It returns
the unconsumed suffix if the prescribed encoding can be skipped, and no result
otherwise.
\end{definition}

\begin{lemma}[Two-tape table reader is a prefix inverse]
\lean{Turing.codeReadTable2_append}
\label{Turing.codeReadTable2_append}
\leanok
\uses{Turing.CodeTable2, Turing.codeReadActionK_append, Turing.codeReadSymbols_append, Turing.codeReadTable2, Turing.codeReadVec_append, Turing.codeTable2Bits}
\difficulty{4}
For every $n\in\mathbb N$, every transition table on states $0,\ldots,n$ and the three
possible scanned symbols (blank, $0$, $1$) on the input tape and on each of two work
tapes, and every bit string $x$, the table reader applied to the serialization of that
table followed by $x$ succeeds and returns the table and $x$. The table serialization
lists the action for each state and each triple of scanned symbols; each action records
the input-head movement, the optional write and movement on each work tape, the optional
output bit, and the optional successor state.
\end{lemma}

\begin{lemma}[Soundness of the two-tape table reader]
\lean{Turing.codeReadTable2_sound}
\label{Turing.codeReadTable2_sound}
\leanok
\uses{Turing.CodeTable2, Turing.codeReadActionK_sound, Turing.codeReadSymbols_sound, Turing.codeReadTable2, Turing.codeReadVec_sound, Turing.codeTable2Bits}
\difficulty{4}
For every $n\in\mathbb N$, bit string $x$, transition table $t$ on states $0,\ldots,n$
and the three possible scanned symbols on the input tape and each of two work tapes, and
suffix $r$, if the table reader with parameter $n$ applied to $x$ returns $(t,r)$, then
$x$ is the serialization of $t$ followed by $r$. The serialization lists the actions in
state and scanned-symbol order, with each action recording the input-head movement, both
work-tape writes and movements, the optional output bit, and the optional successor
state.
\end{lemma}

\begin{lemma}[Erasure law for transition tables]
\lean{Turing.codeEraseTable2}
\label{Turing.codeEraseTable2}
\leanok
\uses{Turing.codeEraseActionK, Turing.codeEraseSymbols, Turing.codeEraseVec, Turing.codeReadActionK, Turing.codeReadSymbols, Turing.codeReadTable2, Turing.codeSkipTable2}
\difficulty{3}
For every natural number $n$ and bit string $x$, decoding a transition table for $n+1$
states, three possible input-symbol values on each of three symbol-indexed dimensions,
and actions of a two-work-tape machine has the same success or failure and leaves the
same suffix as skipping the encoding of that table. In other words, discarding the
decoded table from the reader gives exactly the result of the corresponding table
skipper.
\end{lemma}

\begin{lemma}[Lower bound on table code length]
\lean{Turing.codeTable2Bits_length}
\label{Turing.codeTable2Bits_length}
\leanok
\uses{Turing.Action, Turing.CodeTable2, Turing.actionBitsK, Turing.codeActionK_length, Turing.codeFlatMap_length, Turing.codeTable2Bits}
\difficulty{5}
For every natural number $n$ and every table assigning an action on two work tapes to
each state and each triple of optional input and work symbols over the binary alphabet,
the length of its Boolean code is at least $351(n+1)$.
\end{lemma}

\begin{lemma}[One-tape action decoding]
\lean{Turing.codeReadActionK_one_append}
\label{Turing.codeReadActionK_one_append}
\leanok
\uses{Turing.Action, Turing.actionBits, Turing.actionBits_eq_actionBitsK, Turing.codeReadActionK, Turing.codeReadActionK_append}
\difficulty{3}
For every natural number $n$, every action $a$ of a one-work-tape binary machine with
states indexed by $\{0,\ldots,n\}$, and every bit string $x$, reading the serialization
of $a$ followed by $x$ succeeds and returns $a$ together with $x$.
\end{lemma}

\begin{lemma}[One-tape action reader soundness]
\lean{Turing.codeReadActionK_one_sound}
\label{Turing.codeReadActionK_one_sound}
\leanok
\uses{Turing.Action, Turing.actionBits, Turing.actionBits_eq_actionBitsK, Turing.codeReadActionK, Turing.codeReadActionK_sound}
\difficulty{3}
For every natural number $n$, bit string $x$, action $a$ of a one-work-tape binary
machine with states indexed by $\{0,\ldots,n\}$, and bit string $r$, if reading $x$ as
an action with one work tape and parameter $n$ succeeds with result $(a,r)$, then $x$ is
the serialization of $a$ followed by $r$.
\end{lemma}

\begin{lemma}[One-tape action erasure law]
\lean{Turing.codeEraseActionK_one}
\label{Turing.codeEraseActionK_one}
\leanok
\uses{Turing.codeEraseActionK, Turing.codeReadActionK, Turing.codeSkipAction, Turing.codeSkipAction_eq_codeSkipActionK}
\difficulty{3}
For every natural number $n$ and every bit string $x$, discarding the parsed action and
retaining only the remaining suffix from the one-work-tape action reader over $n+1$
states gives exactly the suffix returned by the transition-record skipper for $n$.
\end{lemma}

\begin{definition}[One-work-tape code format]
\lean{Turing.oneTapeFormat}
\label{Turing.oneTapeFormat}
\leanok
\uses{Turing.Action, Turing.CodeFormat, Turing.CodeTM, Turing.CodeTM.serialize, Turing.actionBits, Turing.actionBits_eq_actionBitsK, Turing.codeActionK_length, Turing.codeEraseActionK_one, Turing.codeEraseSymbols, Turing.codeEraseVec, Turing.codeFallback, Turing.codeFlatMap_length, Turing.codeReadActionK, Turing.codeReadActionK_one_append, Turing.codeReadActionK_one_sound, Turing.codeReadSymbols, Turing.codeReadSymbols_append, Turing.codeReadSymbols_sound, Turing.codeReadVec, Turing.codeReadVec_append, Turing.codeReadVec_sound, Turing.codeSkipAction, Turing.codeSkipRepeat}
The one-work-tape format uses machines with Boolean symbols and $n+1$ states, with
tables assigning an action to each state, possibly blank input symbol, and possibly
blank work-tape symbol. It serializes a machine by pairing the bit representation of $n$
with the unary encoding of its initial state followed by its table bits; its table
reader and skipper parse or erase the corresponding entries. The format has guard
constant $81$, and every malformed string denotes the fixed fallback machine; assembling
a machine records the chosen state count, serialization has the stated paired form,
every machine can be assembled, the reader correctly recovers tables and the remaining
suffix, erasing the read table agrees with skipping it, and every table has at least
$81(n+1)$ bits.
\end{definition}
